\documentclass[13pt]{beamer}
% Fonts and languages
\usepackage{fontspec}
\usepackage{polyglossia}
\usepackage{bidi}
\usepackage{listings}
\usepackage{xcolor}

\setmainlanguage[numerals=western]{arabic}
\setotherlanguage{english}
\newfontfamily\arabicfont[Script=Arabic]{Amiri}
\newfontfamily\arabicfontsf[Script=Arabic]{Amiri}
\newfontfamily\arabicfonttt[Script=Arabic]{Courier New}

% RTL support and nicer bullets for Arabic beamer
\setbeamertemplate{itemize items}[circle]
\setbeamertemplate{itemize subitem}[triangle]

% Listings for C# code blocks

\lstdefinelanguage{CSharp}{
  language={[Sharp]C},
  morekeywords={string,bool,void}
}

\lstset{
  language=CSharp,
  basicstyle=\ttfamily\small,
  keywordstyle=\color{blue},
  stringstyle=\color{red},
  commentstyle=\color{green!50!black},
  showstringspaces=false,
  numbers=none,
  frame=single,
  breaklines=true
}



\title{العمليات الخارجية في لغة \textenglish{C\#}}
\author{الأستاذ محمود اغبارية}
\date{}

\begin{document}

\begin{frame}
  \titlepage
\end{frame}

\begin{frame}[fragile]{مقدمة ومصطلحات أساسية}
العمليات الخارجية تقسّم البرنامج إلى مقاطع صغيرة تؤدي مهام محددة:
\begin{itemize}
  \item تنظيم الكود وجعله أسهل للقراءة والصيانة.
  \item إعادة استخدام نفس الكود أكثر من مرة.
\end{itemize}
\vspace{0.5em}
\textbf{انتبه للفرق بين المصطلحات:}
\begin{itemize}
  \item \textbf{تستقبل:} تقرأ مدخلات من المستخدم عبر \textenglish{Console.ReadLine()}.
  \item \textbf{تتلقى:} تأخذ معطيات (بارامترات) بين القوسين في تعريف العملية.
  \item \textbf{تطبع:} تعرض مخرجات على الشاشة عبر \textenglish{Console.WriteLine(...)}.
  \item \textbf{تعيد:} تُرجع قيمة عبر الأمر \textenglish{return}.
\end{itemize}
\end{frame}

% مبنى تعريف العملية
\begin{frame}[fragile]{مبنى تعريف العملية}
كل عملية خارجية تتكوّن من قسمين رئيسيين:
\begin{itemize}
  \item \textbf{تعريف العملية:} السطر الأوّل ويحتوي على \textenglish{public static}، ونمط القيمة المعادة (أو \textenglish{void})، واسم العملية، والمعطيات بين القوسين مثل \textenglish{(int x, int y)}.
  \item \textbf{جسم العملية:} الأوامر المكتوبة بين القوسين \textenglish{\{ \}} وتُنفّذ عند استدعاء العملية.
\end{itemize}
\end{frame}

% مثال: PrintHello
\begin{frame}[fragile]{مثال $1$: عملية لا تتلقى ولا تعيد}
\begin{english}
\begin{lstlisting}[language=CSharp]
public static void PrintHello()
{
    Console.WriteLine("Hello");
}
\end{lstlisting}
\end{english}
\vspace{0.5em}
\begin{itemize}
  \item \textenglish{public static} $\leftarrow$ جزء ثابت في هذه المرحلة.
  \item \textenglish{void} $\leftarrow$ العملية لا تعيد أي قيمة.
  \item \textenglish{()} $\leftarrow$ الأقواس الفارغة تعني أنّها لا تتلقى معطيات.
\end{itemize}
\end{frame}

% مثال: AddNumbers
\begin{frame}[fragile]{مثال $2$: عملية تتلقى وتعيد}
\begin{english}
\begin{lstlisting}[language=CSharp]
public static int AddNumbers(int a, int b)
{
    int sum = a + b;
    return sum;
}
\end{lstlisting}
\end{english}
\vspace{0.5em}
\begin{itemize}
  \item تعيد قيمة من نمط \textenglish{int}، ويجب أن يطابق نمط القيمة بعد \textenglish{return}.
  \item تتلقى معطيين من نمط \textenglish{int} هما \textenglish{(int a, int b)}.
\end{itemize}
\end{frame}

% استدعاء العملية
\begin{frame}[fragile]{استدعاء العملية في \textenglish{Main}}
\begin{english}
\begin{lstlisting}[language=CSharp]
public static void Main(string[] args)
{
    PrintHello();
    int sum1 = AddNumbers(5, 10);
    int a = 5;
    int b = 10;
    int sum2 = AddNumbers(a, b);
}
\end{lstlisting}
\end{english}
\vspace{0.5em}
\begin{itemize}
  \item \textenglish{PrintHello()} لا تعيد قيمة (\textenglish{void})، لذلك نكتب اسمها فقط.
  \item \textenglish{AddNumbers} تعيد قيمة، فنحفظها في متغير من نفس النمط أو نطبعها.
  \item ترتيب وأنماط المعطيات الممرّرة يجب أن يطابق تعريف العملية.
\end{itemize}
\end{frame}

% استدعاء من فئة أخرى
\begin{frame}[fragile]{استدعاء عملية من فئة أخرى}
نكتب اسم الفئة ثم نقطة ثم اسم العملية:

\begin{english}
\begin{itemize}
  \item Console.WriteLine("Hello World!");
  \item Console.ReadLine();
  \item int.Parse(str); ~ double.Parse(str);
  \item Math.Max(5, 10); ~ Math.Min(5, 10);
  \item Math.Abs(-5); ~ Math.Pow(2, 3);
  \item Math.Sqrt(25); ~ Math.Round(2.5);
\end{itemize}
\end{english}
\end{frame}

\begin{frame}[fragile]{أقسام العمليات الخارجية}
\begin{enumerate}
  \item عمليات لا تتلقى ولا تعيد (تستقبل وتطبع فقط)
  \item عمليات تتلقى ولا تعيد (قسم منها يستقبل أيضًا)
  \item عمليات لا تتلقى وتعيد أعدادًا صحيحة أو عشرية
  \item عمليات تتلقى وتعيد أعدادًا صحيحة أو عشرية
  \item عمليات تعيد \textenglish{string} أو \textenglish{char}
  \item عمليات تعيد \textenglish{bool} وكيفية استخدامها داخل جملة \textenglish{if}
  \item عمليات تستدعي عمليات أخرى
\end{enumerate}
\end{frame}

% القسم الأول
\begin{frame}[fragile]{$1$. عمليات لا تتلقى ولا تعيد (تستقبل وتطبع)}
\begin{english}
\begin{lstlisting}[language=CSharp]
public static void ReadAndPrintSum()
{
    int a = int.Parse(Console.ReadLine());
    int b = int.Parse(Console.ReadLine());
    Console.WriteLine("Sum = " + (a + b));
}

public static void Main()
{
    ReadAndPrintSum();
}
\end{lstlisting}
\end{english}
\end{frame}

% القسم الثاني
\begin{frame}[fragile]{$2$. عمليات تتلقى ولا تعيد}
\begin{english}
\begin{lstlisting}[language=CSharp]
public static void PrintTotalCost(double price)
{
    int quantity = int.Parse(Console.ReadLine());
    double total = price * quantity;
    Console.WriteLine("Total = " + total);
}

public static void Main()
{
    PrintTotalCost(15.5);
}
\end{lstlisting}
\end{english}
\vspace{0.3em}
تتلقى السعر كمعطى، وتستقبل الكمية من المستخدم، ثم تطبع التكلفة.
\end{frame}

% القسم الثالث
\begin{frame}[fragile]{$3$. عمليات لا تتلقى وتعيد عددًا (\textenglish{int / double})}
\begin{english}
\begin{lstlisting}[language=CSharp]
public static int ReadNumber()
{
    int num = int.Parse(Console.ReadLine());
    return num;
}

public static void Main()
{
    int x = ReadNumber();
    Console.WriteLine("The number: " + x);
}
\end{lstlisting}
\end{english}
\end{frame}

% القسم الرابع
\begin{frame}[fragile]{$4$. عمليات تتلقى وتعيد عددًا (\textenglish{int / double})}
\begin{english}
\begin{lstlisting}[language=CSharp]
public static double GetAverage(double a, double b)
{
    double avg = (a + b) / 2.0;
    return avg;
}

public static void Main()
{
    double result = GetAverage(7.5, 9.0);
    Console.WriteLine("Average = " + result);
}
\end{lstlisting}
\end{english}
\end{frame}

% القسم الخامس
\begin{frame}[fragile]{$5$. عمليات تعيد \textenglish{string} أو \textenglish{char}}
\begin{english}
\begin{lstlisting}[language=CSharp]
public static char GetGradeSymbol(int grade)
{
    if (grade >= 55)
    {
        return 'P';
    }
    else
    {
        return 'F';
    }
}

public static void Main()
{
    char sym = GetGradeSymbol(80);
    Console.WriteLine(sym);
}
\end{lstlisting}
\end{english}
\end{frame}

% القسم السادس - شريحة 1
\begin{frame}[fragile]{$6$. تعريف عملية تعيد \textenglish{bool}}
يمكن كتابتها باستعمال \textenglish{if - else} أو بإرجاع الشرط مباشرة:
\begin{english}
\begin{lstlisting}[language=CSharp]
// Way 1: using if - else
public static bool IsEven(int num)
{
    if (num % 2 == 0)
    {
        return true;
    }
    else
    {
        return false;
    }
}

// Way 2: returning the boolean expression directly
public static bool AreEqual(double a, double b)
{
    return (a == b);
}
\end{lstlisting}
\end{english}
\end{frame}

% القسم السادس - شريحة 2
\begin{frame}[fragile]{$6$. استدعاء عملية تعيد \textenglish{bool} داخل جملة \textenglish{if}}
\begin{english}
\begin{lstlisting}[language=CSharp]
public static void Main()
{
    int n = int.Parse(Console.ReadLine());
    if (IsEven(n))
    {
        Console.WriteLine("Even");
    }
    else
    {
        Console.WriteLine("Odd");
    }
}
\end{lstlisting}
\end{english}
\vspace{0.3em}
بما أنّ العملية تعيد \textenglish{true} أو \textenglish{false}، نضع استدعاءها مباشرة بين قوسي \textenglish{if (...)}.
\end{frame}

% القسم السابع
\begin{frame}[fragile]{$7$. عمليات تستدعي عمليات أخرى}
\begin{english}
\begin{lstlisting}[language=CSharp]
public static double AddNumbers(double a, double b)
{
    return a + b;
}

public static double AddFourNumbers(double a, double b, double c, double d)
{
    double sum = AddNumbers(a, b);
    sum = AddNumbers(sum, c);
    sum = AddNumbers(sum, d);
    return sum;
}

public static void Main()
{
    Console.WriteLine(AddFourNumbers(10, 20, 30, 40));
}
\end{lstlisting}
\end{english}
\end{frame}

% ملاحظات هامة
\begin{frame}[fragile]{ملاحظات هامة}
\begin{itemize}
  \item أسماء العمليات تبدأ بحرف كبير وتصف الوظيفة: \textenglish{CalculateSum(), GetMaxValue(), PrintResult()}.
  \item اجعل كل عملية تقوم بمهمة واحدة، وتجنّب تكرار الكود.
  \item يمكن استخدام جمل \textenglish{if - else} داخل العملية، واستدعاء عملية من داخل عملية أخرى.
  \item يمكن تمرير القيمة المعادة من عملية كمعطى لعملية أخرى مباشرة.
\end{itemize}
\end{frame}

\end{document}
